\section{1974: The Visitation and the Declaration of 21 November}

\subsection{The visitation}

In November 1974 two apostolic visitors examined the seminary at Écône.
The visitation itself is not documented here from a Roman act: no decree
of appointment, mandate, or report was located in \work{Acta Apostolicae
Sedis} or on the Holy See's website, and the Society's account is the
only narrative available. That account states that the visitors were
favourably impressed by the seminarians' discipline and studies but
expressed opinions in conversation---about the ordination of married men,
about the mutability of truth, about the Resurrection---that scandalized
the founder.\endnote{``Archbishop Marcel Lefebvre's 1974 Declaration''
(sspx.org, United States District, 21 November 2024), read 2026-07-25,
which prefaces the Declaration with this account of the visitors' remarks
and dates the visit ten days before the Declaration; and the Society's
institutional history cited above. Party sources reporting private
conversations. Bounded negative: no appointment decree, mandate, report,
or communiqué of the 1974 visitation was located on the Holy See website,
and a literal search of the text layer of the scanned \work{Acta
Apostolicae Sedis} 66 (1974) and 67 (1975) returns no occurrence of the
name Marcel Lefebvre at all---the only hits in the 1975 volume belong to
a different person (checked 2026-07-25). What Rome concluded from the visitation is known only from
what followed.}

Whatever was said at Écône, the consequence is dated and undisputed. On
21 November 1974 Archbishop Lefebvre wrote a declaration for his
seminarians. It is the single most consequential text the Society has
ever produced, because the Holy See treated it---not the liturgy, not the
seminary, not the ordinations---as the ground of everything it decided
afterwards.

\subsection{What the Declaration says}

The Declaration is short. It professes adherence to ``Catholic Rome,
Guardian of the Catholic Faith and of the traditions necessary to
preserve this faith,'' and to ``Eternal Rome, Mistress of wisdom and
truth.'' It then refuses ``to follow the Rome of neo-Modernist and
neo-Protestant tendencies which were clearly evident in the Second
Vatican Council and, after the Council, in all the reforms which issued
from it.'' It charges those reforms with contributing ``to the
destruction of the Church, to the ruin of the priesthood, to the
abolition of the Sacrifice of the Mass and of the sacraments.'' It states
that ``no authority, not even the highest in the hierarchy, can force us
to abandon or diminish our Catholic Faith,'' and concludes that the only
faithful attitude is ``a categorical refusal to accept this Reformation,''
while the work of forming priests continues ``without any spirit of
rebellion, bitterness or resentment.''\endnote{Quoted from the Society's
published English text of the Declaration of 21 November 1974 (sspx.org,
21 November 2024), read 2026-07-25. This is a party-hosted translation
whose translator is not identified and whose French original was not
located in an independent witness; the quotations are given as the
Society's own English wording. The Declaration's date is independently
confirmed by Paul VI, who refers to ``pronuntiatio tua, die XXI mensis
Novembris anno MCMLXXIV facta'' (letter of 11 October 1976, section III~B),
and its central phrase is independently attested by the same letter (next
note).}

Rome read it exactly as it was written. In his letter of 11 October 1976
Paul VI reproached Lefebvre for distrusting the Apostolic See itself,
``which you accuse of being a \latin{Rome inclined to neo-modernism and
neo-Protestantism}''---quoting the Declaration's own phrase back at its
author.\endnote{Paul VI, letter of 11 October 1976, section I: ``diffidendi
ipsi Sedi Apostolicae, quam \guillemotleft Romam ad neo-modernismum et
neo-protestantismum proclivem\guillemotright{} esse criminaris.'' The
English words in the text above paraphrase the Latin; no project
translation of the papal letter is offered. This is the strongest
available evidence for the Declaration's wording, because it is a Roman
act quoting the document.} That is the pivot of the whole history: the
Society's founding manifesto asserts a right to judge the acts of the
Roman authority by the criterion of tradition, and the Roman authority
answered that no bishop, alone and without canonical mission, may
determine what tradition contains.

\actrecord{Declaration of Archbishop Marcel Lefebvre, 21 November 1974
(Society's English text; date and key phrase confirmed by Paul VI's
letter of 11 October 1976).}{The founder's public refusal of the
conciliar and post-conciliar reforms as such, and his claim that no
authority may require their acceptance.}{A manifesto written for
seminarians---composed, in the words the Society's own presentation
places in quotation marks, ``in a spirit of doubtlessly excessive
indignation.'' It is evidence of his position in November 1974, not of the Society's institutional practice, and this
account does not adjudicate whether its propositions are compatible with
Catholic doctrine.}

\subsection{Why the Declaration and not the liturgy}

It is often said that the Society was suppressed for celebrating the
older Mass. The documents do not say that. The Cardinals' Commission of
1975, as Paul VI recounted its judgment, concluded that on the foundation
laid by the Declaration no institution or priestly formation answering
the needs of the Church of Christ could be built.\endnote{Paul VI, letter
of 11 October 1976, section III~B: ``sed, quemadmodum merito prorsus
iudicavit Nostra Commissio Cardinalium die VI mensis Maii MCMLXXV, super
tale fundamentum nulla exstrui potest institutio aut formatio sacerdotalis
congruens necessitatibus Ecclesiae Christi.'' The demonstrative
\latin{tale fundamentum} refers to the Declaration named in the preceding
sentence.} The liturgical question is present in the Roman documents, and
Paul VI devoted a long passage of the same letter to the reformed
\work{Ordo Missae} and to the claim that use of the older rite had become
in Lefebvre's hands ``the sign of a false ecclesiology''; but it is
present as a symptom. The stated ground of the 1975 decision is the
Declaration.

That distinction survives every later turn of the story. In 1988 John
Paul II located the root of the schismatic act in ``an incomplete and
contradictory notion of Tradition''; in 2009 Benedict XVI insisted that
the Society's lack of canonical status rested ``not, in the end, \ldots on
disciplinary but on doctrinal reasons''; in 2026 the Dicastery for the
Doctrine of the Faith opened its explanatory note by recalling that
attempts at reconciliation had failed ``from the times of Saint Paul VI
up to the most recent conversations.'' Four pontificates gave the same
diagnosis.
